% ============================================================
%  LinLang — Monografía Doctoral
%  Compilar con: pdflatex monografia.tex  (dos pasadas)
% ============================================================
\documentclass[12pt, a4paper, openany]{report}

% ── Codificación y lengua ────────────────────────────────────────────────────
\usepackage[utf8]{inputenc}
\usepackage[T1]{fontenc}
\usepackage[spanish, es-tabla]{babel}

% ── Tipografía ───────────────────────────────────────────────────────────────
\usepackage{lmodern}
\usepackage{microtype}

% ── Geometría ────────────────────────────────────────────────────────────────
\usepackage[a4paper, top=3cm, bottom=3cm, left=3.5cm, right=2.5cm]{geometry}

% ── Matemáticas ──────────────────────────────────────────────────────────────
\usepackage{amsmath, amssymb, amsthm}
\usepackage{braket}          % notación Dirac |ψ⟩, ⟨φ|
\usepackage{mathtools}
\usepackage{bm}              % vectores en negrita

% ── Gráficos y diagramas ─────────────────────────────────────────────────────
\usepackage{xcolor}
\usepackage{graphicx}
\usepackage{tikz}
\usetikzlibrary{shapes.geometric, arrows.meta, positioning, fit, backgrounds}
\usepackage{float}

% ── Código fuente ────────────────────────────────────────────────────────────
\usepackage{listings}
\usepackage{inconsolata}

\definecolor{codebg}{HTML}{F8F8F8}
\definecolor{codecomment}{HTML}{6A9153}
\definecolor{codekw}{HTML}{0000CD}
\definecolor{codestr}{HTML}{A31515}
\definecolor{codenumber}{HTML}{098658}

\lstdefinelanguage{LinLang}{
  morekeywords={space, let, transform, when, persist, drop, query, query_one,
                hilbert, gate, row, ket, apply, tensor, measure, density,
                partial_trace, braket, purity, kron, trace_alice,
                project, onto, add, scale, dot, norm, print, include},
  sensitive=true,
  morecomment=[l]{\#},
  morestring=[b]{"},
}

\lstset{
  language=LinLang,
  backgroundcolor=\color{codebg},
  basicstyle=\small\ttfamily,
  keywordstyle=\color{codekw}\bfseries,
  commentstyle=\color{codecomment}\itshape,
  stringstyle=\color{codestr},
  numberstyle=\tiny\color{gray},
  numbers=left,
  numbersep=8pt,
  frame=single,
  framesep=4pt,
  rulecolor=\color{gray!40},
  breaklines=true,
  captionpos=b,
  showstringspaces=false,
  tabsize=4,
  xleftmargin=16pt,
}

\lstdefinestyle{go}{
  language=Go,
  backgroundcolor=\color{codebg},
  basicstyle=\small\ttfamily,
  keywordstyle=\color{codekw}\bfseries,
  commentstyle=\color{codecomment}\itshape,
  stringstyle=\color{codestr},
  numbers=left, numberstyle=\tiny\color{gray}, numbersep=8pt,
  frame=single, rulecolor=\color{gray!40},
  breaklines=true, captionpos=b, showstringspaces=false,
  xleftmargin=16pt,
}

% ── Tablas ───────────────────────────────────────────────────────────────────
\usepackage{booktabs}
\usepackage{tabularx}
\usepackage{multirow}
\usepackage{array}
\newcolumntype{L}[1]{>{\raggedright\arraybackslash}p{#1}}
\newcolumntype{C}[1]{>{\centering\arraybackslash}p{#1}}

% ── Entornos matemáticos ─────────────────────────────────────────────────────
\theoremstyle{definition}
\newtheorem{definition}{Definición}[chapter]
\newtheorem{theorem}{Teorema}[chapter]
\newtheorem{proposition}{Proposición}[chapter]
\newtheorem{example}{Ejemplo}[chapter]
\newtheorem{remark}{Observación}[chapter]

% ── Hipervínculos ─────────────────────────────────────────────────────────────
\usepackage[hidelinks, colorlinks=true,
  linkcolor=blue!70!black,
  citecolor=green!50!black,
  urlcolor=blue!70!black]{hyperref}

% ── Bibliografía ─────────────────────────────────────────────────────────────
\usepackage[numbers, sort&compress]{natbib}

% ── Cabeceras ────────────────────────────────────────────────────────────────
\usepackage{fancyhdr}
\pagestyle{fancy}
\fancyhf{}
\fancyhead[L]{\small\leftmark}
\fancyhead[R]{\small LinLang}
\fancyfoot[C]{\thepage}
\renewcommand{\headrulewidth}{0.4pt}

% ── Notación propia ──────────────────────────────────────────────────────────
\newcommand{\linlang}{\textbf{LinLang}}
\newcommand{\linlangQ}{\textbf{LinLang/Q}}
\newcommand{\Hilbert}{\mathcal{H}}
\newcommand{\ket}[1]{\left|#1\right\rangle}
\newcommand{\bra}[1]{\left\langle#1\right|}
\newcommand{\braketnotation}[2]{\left\langle#1\middle|#2\right\rangle}
\newcommand{\density}{\hat{\rho}}
\newcommand{\tr}{\mathrm{Tr}}
\newcommand{\pure}{\mathrm{pure}}

% ============================================================
\begin{document}
% ============================================================

% ── Portada ───────────────────────────────────────────────────────────────────
\begin{titlepage}
  \centering
  \vspace*{2cm}
  {\Large\bfseries Universidad Doctoral en Ciencias de la Computación\par}
  \vspace{1.5cm}
  \rule{\linewidth}{0.5pt}\\[0.4cm]
  {\Huge\bfseries LinLang\par}
  \vspace{0.3cm}
  {\Large Un Lenguaje Algebraico Orientado a Subespacios Vectoriales\\
         con Extensión de Computación Cuántica Distribuida\par}
  \rule{\linewidth}{0.5pt}\\[1.5cm]
  {\large\itshape Monografía Doctoral\par}
  \vspace{0.8cm}
  {\large Aaron Soria\par}
  \vfill
  {\large Mayo 2026\par}
\end{titlepage}

\cleardoublepage

% ── Resumen ───────────────────────────────────────────────────────────────────
\chapter*{Resumen}
\addcontentsline{toc}{chapter}{Resumen}

\linlang{} es un lenguaje de programación experimental basado en los fundamentos del álgebra lineal, en el que cada dominio de datos se modela como un espacio vectorial, cada entidad como un vector y cada relación como una transformación algebraica entre subespacios. El paradigma propuesto ofrece una alternativa declarativa y matemáticamente coherente a los paradigmas orientados a objetos y funcionales.

La presente monografía documenta el diseño, implementación y extensiones del lenguaje en tres dimensiones: \textbf{(1)} el núcleo algebraico clásico (operaciones vectoriales, transformaciones, condicionales); \textbf{(2)} la extensión cuántica \linlangQ{}, que incorpora espacios de Hilbert, estados cuánticos puros y mixtos, puertas unitarias, producto tensorial y medición según la regla de Born; y \textbf{(3)} la capa de persistencia transparente, con soporte nativo de SQLite y Redis accesibles mediante las primitivas \texttt{persist}, \texttt{drop} y \texttt{query} directamente en los archivos \texttt{.lin}.

Se documenta también un backend REST completo (\emph{TODO server}) como caso de estudio, completamente dirigido por el runtime de LinLang, así como la orquestación mediante Docker~Compose y una hoja de ruta hacia la distribución transparente de qubits entre múltiples nodos de cómputo.

\bigskip\noindent\textbf{Palabras clave:} álgebra lineal, programación algebraica, espacios vectoriales, computación cuántica, espacios de Hilbert, sistemas distribuidos, microservicios, Docker, Go.

\chapter*{Abstract}
\addcontentsline{toc}{chapter}{Abstract}

\linlang{} is an experimental programming language grounded in linear algebra, where each data domain is modeled as a vector space, each entity as a vector, and each relationship as an algebraic transformation between subspaces. The proposed paradigm offers a declarative, mathematically coherent alternative to object-oriented and functional programming.

This monograph documents the design, implementation and extensions of the language across three dimensions: \textbf{(1)} the classical algebraic core (vector operations, transforms, conditionals); \textbf{(2)} the quantum extension \linlangQ{}, which incorporates Hilbert spaces, pure and mixed quantum states, unitary gates, tensor product, and Born rule measurement; and \textbf{(3)} a transparent persistence layer with native SQLite and Redis support via the \texttt{persist}, \texttt{drop} and \texttt{query} primitives directly in \texttt{.lin} programs.

A complete REST backend (TODO server) driven entirely by the LinLang runtime is presented as a case study, together with Docker~Compose orchestration and a roadmap toward transparent qubit distribution across heterogeneous compute nodes.

\bigskip\noindent\textbf{Keywords:} linear algebra, algebraic programming, vector spaces, quantum computing, Hilbert spaces, distributed systems, microservices, Docker, Go.

\cleardoublepage

% ── Tabla de contenidos ───────────────────────────────────────────────────────
\tableofcontents
\cleardoublepage

% ============================================================
\chapter{Introducción}
\label{ch:intro}
% ============================================================

La presente monografía expone el diseño e implementación de \linlang{}, un lenguaje de programación basado en los principios del álgebra lineal y orientado a subespacios vectoriales. El objetivo principal es proporcionar un modelo computacional donde las entidades y relaciones de un dominio puedan representarse como subespacios y transformaciones vectoriales, promoviendo una programación algebraica natural, distribuida y escalable.

Desde su concepción, \linlang{} ha evolucionado en tres etapas bien diferenciadas:

\begin{enumerate}
  \item \textbf{Núcleo algebraico clásico}: definición de espacios vectoriales (\texttt{space}), instanciación de vectores (\texttt{let}), transformaciones entre espacios (\texttt{transform}) y condicionales algebraicos (\texttt{when}).
  \item \textbf{Extensión cuántica \linlangQ{}}: introducción de espacios de Hilbert sobre $\mathbb{C}^n$, estados cuánticos puros, matrices de densidad, puertas unitarias, producto de Kronecker y medición probabilística según la regla de Born.
  \item \textbf{Capa de persistencia y distribución}: backend de base de datos transparente (SQLite~/ Redis) con primitivas nativas en el lenguaje, orquestación mediante Docker~Compose y una hoja de ruta hacia la distribución de qubits entre nodos heterogéneos.
\end{enumerate}

\section{Contribuciones}

Las principales contribuciones de este trabajo son:

\begin{itemize}
  \item Un paradigma de programación algebraico donde los tipos son espacios vectoriales y las funciones son transformaciones lineales entre ellos.
  \item La unificación del paradigma clásico y cuántico bajo el mismo formalismo: los subespacios de Hilbert son una instancia del concepto de espacio vectorial ya presente en el núcleo del lenguaje.
  \item Primitivas de persistencia de primera clase (\texttt{persist}, \texttt{drop}, \texttt{query}) que generan el esquema de base de datos directamente desde la declaración del espacio vectorial.
  \item Un prototipo funcional en Go con parser, intérprete, extensión cuántica y servidor HTTP de demostración, disponible en \url{https://github.com/Bulldrill/Linglang}.
\end{itemize}

\section{Estructura del documento}

El capítulo~\ref{ch:motivacion} presenta la motivación y los fundamentos teóricos. El capítulo~\ref{ch:diseno} describe el diseño del lenguaje. El capítulo~\ref{ch:implementacion} detalla la arquitectura de la implementación. El capítulo~\ref{ch:quantum} desarrolla la extensión cuántica. El capítulo~\ref{ch:persistencia} describe la capa de persistencia transparente. El capítulo~\ref{ch:distribucion} analiza la orquestación y la distribución. El capítulo~\ref{ch:resultados} presenta los resultados. El capítulo~\ref{ch:futuro} establece la hoja de ruta, y el capítulo~\ref{ch:conclusiones} cierra con las conclusiones.

% ============================================================
\chapter{Motivación y Fundamentos Teóricos}
\label{ch:motivacion}
% ============================================================

\section{Motivación}

La programación tradicional impone estructuras ajenas a la naturaleza matemática de los datos. Los sistemas orientados a objetos fragmentan la lógica del dominio en jerarquías de clases, mientras que los sistemas funcionales requieren codificar explícitamente la transformación de cada entidad. LinLang propone un enfoque alternativo fundamentado en cuatro principios:

\begin{definition}[Principio de Espacialidad]
Todo dominio de datos $D$ se modela como un espacio vectorial $V_D \subseteq \mathbb{R}^n$, cuyas dimensiones corresponden a los atributos del dominio.
\end{definition}

\begin{definition}[Principio de Vectorialidad]
Toda entidad $e \in D$ es un vector $\bm{v} \in V_D$ cuyos componentes representan los valores de los atributos.
\end{definition}

\begin{definition}[Principio de Transformabilidad]
Toda relación entre dominios $D_1$ y $D_2$ se expresa como una transformación $T: V_{D_1} \times V_{D_2} \to V_{D_3}$.
\end{definition}

\begin{definition}[Principio de Condicionalidad Algebraica]
Toda condición sobre una entidad es una evaluación escalar sobre una componente del vector correspondiente.
\end{definition}

\section{Fundamentos de Álgebra Lineal}

Sea $\mathbb{F}$ un cuerpo ($\mathbb{R}$ o $\mathbb{C}$). Un \emph{espacio vectorial} sobre $\mathbb{F}$ es un conjunto $V$ dotado de dos operaciones cerradas: suma vectorial y multiplicación escalar, que satisfacen los axiomas estándar de asociatividad, conmutatividad, existencia de elemento neutro e inverso.

En \linlang{}, la declaración

\begin{lstlisting}
space Tarea: id: Real, prioridad: Real, estado: Real, categoria: Real
\end{lstlisting}

define formalmente el espacio $V_{\text{Tarea}} = \mathbb{R}^4$ con base canónica $\{e_{\text{id}},\ e_{\text{prioridad}},\ e_{\text{estado}},\ e_{\text{categoria}}\}$.

\section{Fundamentos de Mecánica Cuántica}

La extensión \linlangQ{} se asienta sobre la formulación matemática de la mecánica cuántica en términos de espacios de Hilbert.

\begin{definition}[Espacio de Hilbert]
Un \emph{espacio de Hilbert} $\Hilbert$ es un espacio vectorial sobre $\mathbb{C}$ dotado de un producto interno $\braket{\cdot|\cdot}: \Hilbert \times \Hilbert \to \mathbb{C}$ completo bajo la norma inducida $\|\psi\| = \sqrt{\braketnotation{\psi}{\psi}}$.
\end{definition}

\begin{definition}[Estado cuántico puro]
Un \emph{estado puro} es un vector normalizado $\ket{\psi} \in \Hilbert$ con $\|\ket{\psi}\| = 1$. En una base computacional $\{\ket{0}, \ket{1}, \ldots, \ket{n-1}\}$:
\[
  \ket{\psi} = \sum_{i=0}^{n-1} \alpha_i \ket{i},\quad \sum_{i=0}^{n-1} |\alpha_i|^2 = 1.
\]
\end{definition}

\begin{definition}[Matriz de densidad]
Un \emph{estado mixto} se describe mediante un operador densidad $\density = \sum_i p_i \ket{\psi_i}\bra{\psi_i}$, donde $p_i \geq 0$ y $\sum_i p_i = 1$. La \emph{pureza} del estado es $\tr(\density^2) \in [1/n,\, 1]$, igual a $1$ si y solo si el estado es puro.
\end{definition}

\begin{definition}[Puerta cuántica]
Una \emph{puerta cuántica} es un operador unitario $U: \Hilbert \to \Hilbert$ tal que $U^\dagger U = I$.
\end{definition}

\begin{definition}[Regla de Born]
La probabilidad de medir el estado base $\ket{k}$ al observar $\ket{\psi}$ es $P(k) = |\braketnotation{k}{\psi}|^2 = |\alpha_k|^2$.
\end{definition}

\section{Comparativa con Otros Paradigmas}

\begin{table}[H]
\centering
\caption{Comparativa de LinLang con paradigmas existentes}
\label{tab:comparativa}
\begin{tabularx}{\linewidth}{@{}L{2.8cm}L{4cm}L{5.5cm}@{}}
\toprule
\textbf{Paradigma} & \textbf{Enfoque} & \textbf{Relación con LinLang} \\
\midrule
OOP & Clases, herencia, encapsulación & LinLang abstrae entidades como vectores; no hay jerarquías de clases \\
\addlinespace
Funcional & Composición de funciones puras & Los transforms de LinLang son funciones lineales entre espacios; la composición es el producto de transformaciones \\
\addlinespace
Data Spatial (DSP)\cite{mars2025} & Topología de datos con nodos y walkers & LinLang simplifica a álgebra lineal pura sin walkers ni topologías de grafos \\
\addlinespace
Declarativo (SQL) & Consultas sobre relaciones & \texttt{query Space where dim op val} es el análogo algebraico de \texttt{SELECT WHERE} \\
\addlinespace
Q\#, Qiskit, Cirq & Lenguajes cuánticos especializados & LinLang/Q unifica el cómputo clásico y cuántico bajo el mismo formalismo de espacios vectoriales \\
\addlinespace
OpenQASM & Ensamblador cuántico & LinLang/Q opera a nivel de abstracción más alto; el compilador genera el circuito bajo demanda \\
\bottomrule
\end{tabularx}
\end{table}

\subsection{Comparativa Detallada con Lenguajes Cuánticos Establecidos}
\label{sec:comparativa-detallada}

La fila de la tabla anterior dedicada a Qiskit/Q\#/Cirq/OpenQASM resume la relación en una frase; esta sección la desarrolla lenguaje por lenguaje, sobre cuatro ejes: modelo de ejecución, sistema de tipos, soporte nativo de distribución, e integración con cómputo clásico.

\begin{table}[H]
\centering
\caption{\linlangQ{} frente a Qiskit, Q\#, Cirq y OpenQASM}
\label{tab:comparativa-detallada}
\begin{tabularx}{\linewidth}{@{}L{1.8cm}L{3.3cm}L{3.3cm}L{2.3cm}L{3.3cm}@{}}
\toprule
\textbf{Lenguaje} & \textbf{Modelo de ejecución} & \textbf{Sistema de tipos} & \textbf{Distribución nativa} & \textbf{Integración clásica} \\
\midrule
Qiskit & API Python que construye un objeto \texttt{QuantumCircuit}, transpilado y enviado a un backend (simulador o IBM Quantum) & Dinámico (Python); los circuitos son datos en tiempo de ejecución, no un tipo del lenguaje & No — el usuario orquesta multi-backend manualmente con el SDK & Total (es una librería Python); sin relación algebraica con el cómputo clásico \\
\addlinespace
Q\# & Lenguaje dedicado, compilado; \texttt{operation}/\texttt{function} con \texttt{Qubit} como tipo primitivo gestionado por el runtime (\texttt{use}/\texttt{borrow}) & Estático, con \texttt{Qubit} opaco (no se puede copiar ni inspeccionar su amplitud) & No — el modelo asume un único QPU lógico; el runtime de ejecución puede particionar, pero no es parte del lenguaje & Fuerte (host .NET/Python), pero el código cuántico y clásico son lenguajes/tipos distintos \\
\addlinespace
Cirq & API Python centrada en \texttt{Circuit}/\texttt{Moment}, con foco en NISQ y control fino de la disposición temporal de puertas & Dinámico (Python), igual que Qiskit & No — orientado a un único dispositivo o simulador por ejecución & Total (librería Python), misma limitación que Qiskit \\
\addlinespace
OpenQASM & Lenguaje de ensamblador: lista plana de instrucciones de puertas sobre registros de qubits/bits clásicos & Estático pero mínimo (\texttt{qubit}, \texttt{bit}); sin abstracción algebraica & No — es el formato de intercambio, no una capa de orquestación & Ninguna; es el objetivo de compilación de Qiskit/Cirq/Q\#, no un lenguaje de autoría \\
\addlinespace
\textbf{\linlangQ{}} & El mismo \texttt{Runtime} que ejecuta el núcleo clásico; \texttt{apply}/\texttt{measure} se resuelven contra un \texttt{QuantumBackend} intercambiable (simulador, GPU, IBM, IonQ) & Estático a nivel de espacio: \texttt{HilbertSpace} es una instancia de \texttt{Space} sobre $\mathbb{C}$, con la misma sintaxis \texttt{space}/\texttt{let} que los vectores clásicos & \textbf{Sí} — \texttt{QuantumChannel}, \texttt{Teleport} y \texttt{ScheduleCircuit} son primitivas del lenguaje/runtime, no una capa externa (\S\ref{sec:hilbert-distribuido}) & Total por construcción: no hay dos lenguajes — \texttt{space}, \texttt{let}, \texttt{transform} son los mismos para $\mathbb{R}^n$ y $\mathbb{C}^n$ \\
\bottomrule
\end{tabularx}
\end{table}

Dos diferencias merecen explicitarse más allá de la tabla:

\begin{itemize}
  \item \textbf{Unificación clásico-cuántica.} Qiskit, Cirq y Q\# tratan lo cuántico como un subsistema que un programa \emph{clásico} invoca (una librería o un lenguaje huésped distinto). \linlangQ{} no introduce un segundo lenguaje: \texttt{hilbert} es a \texttt{space} lo que $\mathbb{C}^n$ es a $\mathbb{R}^n$, y el mismo \texttt{Runtime} interpreta ambos. La contribución no es un nuevo lenguaje cuántico, sino la demostración de que la extensión es innecesaria si el núcleo ya es álgebra lineal.
  \item \textbf{Distribución como primitiva, no como infraestructura externa.} Ninguno de los cuatro lenguajes comparados ofrece una primitiva de \emph{migración} de un qubit entre dispositivos — la pregunta no se plantea porque todos asumen un único QPU lógico por ejecución. \texttt{teleport(psi)} y \texttt{ScheduleCircuit} (\S\ref{sec:hilbert-distribuido}, issues \#14, \#18) convierten esa migración en una llamada de runtime, en la misma posición sintáctica que \texttt{persist} convierte el almacenamiento en una llamada de runtime (Capítulo~\ref{ch:persistencia}).
\end{itemize}

% ============================================================
\chapter{Diseño del Lenguaje}
\label{ch:diseno}
% ============================================================

\section{Principios de Diseño}

LinLang sigue tres principios rectores de diseño:

\textbf{Transparencia algebraica}: la sintaxis debe reflejar directamente los objetos matemáticos subyacentes. Un programador con conocimientos de álgebra lineal puede leer un programa LinLang como si leyera una especificación matemática.

\textbf{Distribución transparente}: el mismo archivo \texttt{.lin} debe ejecutarse sin modificación independientemente del backend de cómputo (local, contenedor Docker, pod Kubernetes, QPU simulado o real).

\textbf{Extensibilidad vertical}: el núcleo clásico y la extensión cuántica deben ser instancias del mismo formalismo, no capas separadas. Un espacio de Hilbert es un espacio vectorial sobre $\mathbb{C}$, por lo que \texttt{hilbert} es una especialización de \texttt{space}.

\section{Núcleo Clásico}

\subsection{Espacios y Vectores}

\begin{lstlisting}[caption={Definición de espacios y vectores}]
# Definición de espacios vectoriales (dominios)
space Personas:   edad: Real, ingresos: Real, energia: Real
space Mascotas:   edad: Real, especie: Real
space Adopciones: persona_id: Real, mascota_id: Real, compat: Real

# Instanciación de vectores (entidades)
let p1 = Personas[25, 40000, 0.8]
let m1 = Mascotas[3, 1]
\end{lstlisting}

\subsection{Transformaciones}

Las transformaciones son relaciones algebraicas entre espacios. La función $f: V_{D_1} \times V_{D_2} \to V_{D_3}$ se expresa mediante un cuerpo declarativo donde cada dimensión del codominio es una expresión aritmética sobre las entradas:

\begin{lstlisting}[caption={Transform declarativo}]
transform adoptar: Personas x Mascotas -> Adopciones {
    persona_id = p.id
    mascota_id = m.id
    compat     = p.energia * m.edad / 10
}
\end{lstlisting}

El evaluador de expresiones soporta $+$, $-$, $\times$, $/$ con acceso a componentes mediante \texttt{p.dim} y \texttt{m.dim}.

\subsection{Operaciones Vectoriales Nativas}

\begin{lstlisting}[caption={Operaciones sobre vectores}]
let suma      = add(p1, p2)         # suma vectorial
let escalado  = scale(p1, 2.0)      # multiplicacion escalar
let similitud = dot(p1, p2)         # producto punto (escalar)
let magnitud  = norm(p1)            # norma euclidiana
let perfil    = project p1 onto Perfil  # proyeccion a subespacio
\end{lstlisting}

\subsection{Condicionales Algebraicos}

\begin{lstlisting}[caption={Condicionales sobre componentes vectoriales}]
when a1.compat >= 0.2:
    approve(a1)

when a1.compat < 0.1:
    reject(a1)
\end{lstlisting}

Los operadores soportados son $>, <, \geq, \leq, ==$, $\neq$.

\section{Extensión Cuántica (LinLang/Q)}

\subsection{Espacios de Hilbert y Estados}

\begin{lstlisting}[caption={Declaración cuántica en LinLang/Q}]
hilbert Qubit:  dim 2
hilbert QReg2:  dim 4

# Estados base
let ket0 = ket(Qubit, 1, 0)          # |0>
let ket1 = ket(Qubit, 0, 1)          # |1>
let plus = ket(Qubit, 0.7071, 0.7071) # |+>
\end{lstlisting}

\subsection{Puertas Cuánticas}

Las puertas se declaran como matrices con verificación automática de unitariedad ($U^\dagger U = I$):

\begin{lstlisting}[caption={Declaración de puertas unitarias}]
gate H: Qubit -> Qubit {
    row [0.7071067811865476,  0.7071067811865476]
    row [0.7071067811865476, -0.7071067811865476]
}
gate CNOT: QReg2 -> QReg2 {
    row [1, 0, 0, 0]
    row [0, 1, 0, 0]
    row [0, 0, 0, 1]
    row [0, 0, 1, 0]
}
\end{lstlisting}

\subsection{Operaciones Cuánticas}

\begin{lstlisting}[caption={Operaciones nativas de LinLang/Q}]
let q_plus   = apply(H, ket0)          # U|psi>
let q00      = tensor(ket0, ket0)      # |psi> otimes |phi>
let rho      = density(q_plus)         # rho = |psi><psi|
let rho_A    = trace_alice(rho, 2)     # traza parcial
let overlap  = braket(ket0, plus)      # <phi|psi>
let pur      = purity(rho)             # Tr(rho^2)
let result   = measure(q_plus)         # colapso (regla de Born)
\end{lstlisting}

\subsection{Condicionales Cuánticos}

\begin{lstlisting}[caption={Condicionales sobre propiedades cuánticas}]
when purity(rho) >= 0.99:
    estado_puro(rho)

when purity(rho_A) < 0.6:
    entrelazado(rho_A)

when qnorm(plus) >= 0.99:
    normalizado(plus)
\end{lstlisting}

\section{Primitivas de Persistencia}

\begin{lstlisting}[caption={Persistencia nativa de primera clase}]
# INSERT / UPDATE
persist t1

# SELECT con filtro algebraico
let todas    = query Tarea
let urgentes = query Tarea where prioridad == 3
let t        = query_one Tarea where id == 5

# UPDATE = transform + persist
let t_done = completar(t, t)
persist t_done

# DELETE
drop t
\end{lstlisting}

El backend se selecciona mediante la variable de entorno \texttt{LINLANG\_DB}:
\begin{itemize}
  \item \texttt{sqlite://linlang.db} --- SQLite embebido.
  \item \texttt{redis://localhost:6379/0} --- Redis.
  \item \texttt{memory://} --- memoria (tests).
\end{itemize}

% ============================================================
\chapter{Implementación}
\label{ch:implementacion}
% ============================================================

\section{Arquitectura del Prototipo}

El prototipo está implementado en Go~1.22 y organizado en los siguientes paquetes:

\begin{figure}[H]
\centering
\begin{tikzpicture}[
  box/.style={rectangle, draw=gray!70, fill=gray!10, rounded corners=3pt,
              minimum width=2.8cm, minimum height=0.8cm, font=\small\ttfamily},
  arrow/.style={-Latex, gray!70},
  every node/.style={font=\small}
]
  \node[box, fill=blue!10]  (main)   at (0,4)    {main.go};
  \node[box, fill=blue!10]  (server) at (4,4)    {cmd/todo-server};
  \node[box, fill=green!10] (parser) at (0,2.5)  {parser/};
  \node[box, fill=green!10] (db)     at (3,2.5)  {parser/db.go};
  \node[box, fill=green!10] (qp)     at (6,2.5)  {parser/quantum.go};
  \node[box, fill=yellow!20](core)   at (3,1)    {core/};
  \node[box, fill=orange!15](store)  at (0,1)    {store/};
  \node[draw=gray!50, dashed, fit=(core)(store)(qp), inner sep=6pt] {};
  \draw[arrow] (main) -- (parser);
  \draw[arrow] (server) -- (parser);
  \draw[arrow] (parser) -- (core);
  \draw[arrow] (db) -- (store);
  \draw[arrow] (parser) -- (db);
  \draw[arrow] (parser) -- (qp);
  \draw[arrow] (qp) -- (core);
\end{tikzpicture}
\caption{Arquitectura de módulos del prototipo LinLang}
\label{fig:arquitectura}
\end{figure}

\section{Paquete \texttt{core}}

El paquete central define las estructuras de datos algebraicas:

\begin{itemize}
  \item \textbf{\texttt{Space}}: espacio vectorial con nombre y lista de dimensiones.
  \item \textbf{\texttt{Vector}}: instancia con ID auto-incremental y slice de \texttt{float64}.
  \item \textbf{\texttt{Transform}}: función Go encapsulada con dominio y codominio tipados.
  \item \textbf{\texttt{Conditional}}: predicado con clausura \texttt{GetValue~func()~float64}, decoupled del tipo de dato subyacente.
  \item \textbf{\texttt{Expr}}: árbol de expresiones aritméticas con evaluador recursivo descendente.
  \item \textbf{\texttt{HilbertSpace}}, \textbf{\texttt{QuantumState}}, \textbf{\texttt{DensityMatrix}}, \textbf{\texttt{Gate}}: tipos cuánticos sobre $\mathbb{C}$.
\end{itemize}

\subsection{Parser de Expresiones}

El evaluador de expresiones (\texttt{core/expressions.go}) implementa un parser recursivo descendente con precedencia de operadores y soporte para acceso a dimensiones por nombre:

\[
  \text{expr} \to \text{addSub} \to \text{mulDiv} \to \text{unary} \to \text{atom}
\]

donde \texttt{atom} reconoce literales numéricos, accesos \texttt{p.dim}, \texttt{m.dim} y expresiones entre paréntesis. Un punto crítico de diseño es que las dimensiones del espacio vectorial tienen precedencia sobre el identificador interno \texttt{id}:

\begin{lstlisting}[style=go, caption={Prioridad de dimensiones sobre el contador interno}]
// Space dimensions take priority over the internal-counter shorthand.
for i, d := range v.Space.Dimensions {
    if d == e.Dim {
        return v.Values[i]   // dimension named "id" in the space
    }
}
if e.Dim == "id" {
    return float64(v.ID)     // fallback: internal counter
}
\end{lstlisting}

Este invariante evita una clase de errores donde transforms que preservan el identificador del dominio producían vectores con identificadores internos erróneos.

\section{Paquete \texttt{parser}}

El \texttt{Runtime} es la estructura central del intérprete. Mantiene mapas de espacios, vectores, escalares, transforms, colecciones y espacios de Hilbert, así como el estado del backend de persistencia:

\begin{lstlisting}[style=go, caption={Campos del Runtime}]
type Runtime struct {
    Spaces          map[string]*core.Space
    Vectors         map[string]*core.Vector
    Scalars         map[string]float64
    Transforms      map[string]*core.Transform
    Collections     map[string][]*core.Vector
    HilbertSpaces   map[string]*core.HilbertSpace
    QuantumStates   map[string]*core.QuantumState
    DensityMatrices map[string]*core.DensityMatrix
    Gates           map[string]*core.Gate
    Store           store.Backend
    // ... estado de parsing multi-linea
}
\end{lstlisting}

El parser es de un solo paso, línea por línea, con dos excepciones: los cuerpos de \texttt{transform~\{...\}} y \texttt{gate~\{...\}} son multi-línea con un autómata de estado (\texttt{inTransform}, \texttt{inGate}).

Los comentarios inline (\texttt{ \#}) se eliminan en \texttt{ParseLine} antes del despacho, evitando que tokens dentro del texto del comentario contaminen el análisis sintáctico.

\section{Paquete \texttt{store}}

La capa de persistencia implementa el patrón \emph{Repository} con una interfaz agnóstica al driver:

\begin{lstlisting}[style=go, caption={Interfaz Backend}]
type Backend interface {
    Upsert(space *core.Space, values []float64) error
    Query(space *core.Space,
          filter func([]float64) bool) ([]*core.Vector, error)
    Delete(space *core.Space, id float64) error
    Close() error
}
\end{lstlisting}

La implementación SQLite (\texttt{store/sqlite.go}) genera dinámicamente las tablas al primer \texttt{persist}, mapeando la declaración \texttt{space} directamente al esquema relacional: la primera dimensión es \texttt{PRIMARY~KEY}, el resto son \texttt{REAL~NOT~NULL~DEFAULT~0}. La implementación Redis (\texttt{store/redis.go}) almacena cada vector como un Redis Hash bajo la clave \texttt{linlang:\{Space\}:\{id\}}, manteniendo un Set de IDs por espacio para consultas eficientes.

% ============================================================
\chapter{Extensión Cuántica: LinLang/Q}
\label{ch:quantum}
% ============================================================

\section{Motivación}

Los espacios de Hilbert son espacios vectoriales sobre $\mathbb{C}$. La extensión \linlangQ{} es, por tanto, una instancia del núcleo algebraico de LinLang con las siguientes particularidades:

\begin{itemize}
  \item Los vectores son vectores de amplitudes complejas normalizados.
  \item Las transformaciones son puertas unitarias verificadas.
  \item Las condicionales pueden evaluar propiedades de matrices de densidad.
\end{itemize}

Esta coherencia con el núcleo es una contribución de diseño explícita: no existe un intérprete cuántico separado, sino que \texttt{parser/quantum.go} extiende el mismo \texttt{Runtime}.

\section{Modelo de Tipos Cuánticos}

\begin{figure}[H]
\centering
\begin{tikzpicture}[
  node distance=1.4cm,
  box/.style={rectangle, draw=gray!60, rounded corners=3pt,
              minimum width=3.2cm, minimum height=0.7cm,
              fill=blue!8, font=\small\ttfamily},
  arrow/.style={-Latex, blue!50}
]
  \node[box] (hs) {HilbertSpace};
  \node[box, right=of hs] (qs) {QuantumState};
  \node[box, right=of qs] (dm) {DensityMatrix};
  \node[box, below=of hs] (gate) {Gate};
  \node[box, below=of qs] (kron) {KronGate};
  \draw[arrow] (hs) -- node[above,font=\tiny]{tiene dim} (qs);
  \draw[arrow] (qs) -- node[above,font=\tiny]{toDensity()} (dm);
  \draw[arrow] (gate) -- node[right,font=\tiny]{apply()} (qs);
  \draw[arrow] (kron) -- node[above,font=\tiny]{G1$\otimes$G2} (gate);
\end{tikzpicture}
\caption{Jerarquía de tipos en LinLang/Q}
\end{figure}

\section{Estado de Bell y Entrelazamiento}

El estado de Bell $\ket{\Phi^+} = \frac{1}{\sqrt{2}}(\ket{00} + \ket{11})$ se genera en LinLang/Q mediante la composición de las puertas Hadamard y CNOT:

\begin{lstlisting}[caption={Generación del estado de Bell en LinLang/Q}]
let II      = kron(I_Qubit, I_Qubit)
let HII     = kron(H_Qubit, II)         # H otimes I otimes I
let I_CNOT  = kron(I_Qubit, CNOT_QReg2) # I otimes CNOT
let h0      = apply(H, ket0)
let pre_bell = tensor(h0, ket0)
let bell    = apply(CNOT, pre_bell)
\end{lstlisting}

La matriz de densidad del estado de Bell tiene $\tr(\density^2) = 1.0$ (estado puro), mientras que la traza parcial sobre cualquier subsistema produce $\density_A = I/2$ con pureza $0.5$ (mezcla maximal), confirmando el entrelazamiento:

\begin{equation}
  \density_A = \tr_B(\ket{\Phi^+}\bra{\Phi^+}) = \frac{I}{2}
  \implies \tr(\density_A^2) = \frac{1}{2}
\end{equation}

\section{Protocolo de Teleportación Cuántica}

El protocolo de teleportación de Bennett et al.\cite{bennett1993} ha sido implementado íntegramente en LinLang/Q (archivo \texttt{examples/teleportacion.lin}). El protocolo transfiere el estado $\ket{\psi}$ de Alice a Bob usando un par de Bell compartido y dos bits clásicos de comunicación.

Las etapas del protocolo y sus correspondientes instrucciones en LinLang/Q son:

\begin{enumerate}
  \item \textbf{Preparación}: \texttt{let step0 = tensor(psi, q00)}
  \item \textbf{Par de Bell} ($\ket{\Phi^+}$ entre qubits 1 y 2): \texttt{apply(IHI, ...)} $\to$ \texttt{apply(I\_CNOT12, ...)}
  \item \textbf{CNOT de Alice}: \texttt{apply(CNOT01\_I, step2)}
  \item \textbf{Hadamard de Alice}: \texttt{apply(HII, step3)}
  \item \textbf{Medición}: \texttt{measure(step4)} $\to$ bits clásicos $m_0, m_1$
  \item \textbf{Corrección de Bob}: si $m_1=1$ aplica $X$; si $m_0=1$ aplica $Z$
\end{enumerate}

La verificación formal produce:
\begin{itemize}
  \item Todas las probabilidades en \texttt{step4} son $0.125$ (información sobre $\ket{\psi}$ uniformemente distribuida).
  \item $\tr(\density_{\text{Bob}}^2) = 0.5$ antes de la corrección (sin comunicación clásica, Bob no puede conocer $\ket{\psi}$).
  \item Fidelidad $|\braketnotation{\psi}{\psi_{\text{Bob}}}|^2 = 1.0$ después de la corrección.
\end{itemize}

Esta implementación tiene relevancia directa para el capítulo de distribución: la teleportación es exactamente el mecanismo por el cual un qubit ``migra'' entre nodos en la arquitectura distribuida proyectada.

\section{Construcción de Puertas Compuestas}
\label{sec:puertas-compuestas}

La operación \texttt{kron(G1, G2)} implementa el producto de Kronecker de dos puertas, produciendo la puerta compuesta en el espacio producto:

\[
  (G_1 \otimes G_2)_{i \cdot n_2 + k,\; j \cdot n_2 + l} = (G_1)_{ij} \cdot (G_2)_{kl}
\]

Esta operación es esencial para construir puertas de $n$ qubits a partir de puertas de 1 o 2 qubits. Por ejemplo, $H \otimes I \otimes I$ actúa sobre el primer qubit de un registro de 3 qubits:

\begin{lstlisting}
let HI  = kron(H_Qubit, I_Qubit)  # H otimes I sobre QReg2
let HII = kron(H_Qubit, II)       # H otimes I otimes I sobre QReg3
\end{lstlisting}

\section{Formalización Algebraica de \linlangQ{}}
\label{sec:formalizacion-algebra}

Las secciones anteriores presentaron \linlangQ{} operacionalmente: qué hace cada instrucción. Esta sección formaliza \emph{por qué} esas instrucciones son las correctas, situando cada primitiva del lenguaje como la reificación sintáctica de una construcción algebraica estándar de la mecánica cuántica.

\subsection{El producto tensorial como construcción de programación}

En la formulación matemática estándar, los espacios de Hilbert de dimensión finita forman una categoría $\mathbf{Hilb}$ cuyos objetos son espacios $\mathcal{H}^n$ y cuyos morfismos son aplicaciones lineales; el producto tensorial $\otimes: \mathbf{Hilb} \times \mathbf{Hilb} \to \mathbf{Hilb}$ es un funtor bifuntorial que combina dos espacios en uno de dimensión producto. La mayoría de los lenguajes para computación cuántica (Qiskit, Cirq) tratan $\otimes$ como un detalle de representación interno al \emph{registro} de qubits: el programador declara qubits individuales y el framework decide implícitamente cómo se combinan sus espacios.

\linlangQ{} toma una decisión de diseño distinta: \emph{reifica} $\otimes$ como una operación de primera clase del lenguaje.

\begin{definition}[Producto tensorial como operación del runtime]
\label{def:tensor-runtime}
Sean $\mathcal{H}_1, \mathcal{H}_2$ espacios de Hilbert con $\dim \mathcal{H}_1 = n_1$, $\dim \mathcal{H}_2 = n_2$, y $\ket{\psi_1} \in \mathcal{H}_1$, $\ket{\psi_2} \in \mathcal{H}_2$. La instrucción \texttt{tensor(psi1, psi2)} del runtime LinLang computa el vector de amplitudes de $\ket{\psi_1}\otimes\ket{\psi_2} \in \mathcal{H}_1 \otimes \mathcal{H}_2$ ($\dim = n_1 n_2$) vía el producto de Kronecker de las amplitudes, y registra (o reutiliza) el \texttt{HilbertSpace} producto correspondiente en el \texttt{Runtime}. Análogamente, \texttt{kron(G1, G2)} construye el morfismo producto $G_1 \otimes G_2$ sobre el espacio producto.
\end{definition}

La consecuencia de diseño es que un circuito de $n$ qubits en \linlangQ{} \emph{no} es una lista de puertas sobre índices de un registro implícito, sino una expresión algebraica explícita en el espacio producto —tal como se observó en la Sección~\ref{sec:puertas-compuestas} anterior con \texttt{HII = kron(H\_Qubit, II)}— lo cual hace que la correspondencia entre la sintaxis y el álgebra subyacente sea literal, no incidental.

\subsection{Transformaciones unitarias como morfismos entre subespacios de Hilbert}

\begin{definition}[Puerta cuántica como morfismo]
\label{def:gate-morfismo}
Una declaración \texttt{gate G: Dominio -> Codominio \{ ... \}} define un morfismo $G: \mathcal{H}_{\text{dom}} \to \mathcal{H}_{\text{cod}}$ en $\mathbf{Hilb}$, representado por una matriz $U \in \mathbb{C}^{n\times n}$. $G$ pertenece a la subcategoría de morfismos \emph{unitarios} si y solo si preserva el producto interno:
\[
  \forall \ket{\varphi},\ket{\psi} \in \mathcal{H}_{\text{dom}}: \quad \braketnotation{U\varphi}{U\psi} = \braketnotation{\varphi}{\psi} \iff U^\dagger U = I.
\]
\end{definition}

\texttt{Gate.IsUnitary()} (\texttt{core/gates.go}) verifica exactamente esta condición dentro de una tolerancia $10^{-9}$, y \texttt{finalizeGate} (\texttt{parser/quantum.go}) la evalúa automáticamente al cerrar cada declaración \texttt{gate}. La composición de morfismos —aplicar $G_2$ después de $G_1$— corresponde a dos llamadas sucesivas a \texttt{apply}, y la ley de composición categórica (asociatividad) se cumple trivialmente porque \texttt{apply} es multiplicación matricial estándar:
\[
  \texttt{apply(G2, apply(G1, psi))} \;=\; (G_2 G_1)\ket{\psi} \;=\; \texttt{apply(kron\_compose(G2,G1), psi)}.
\]

\begin{proposition}[Cierre de la subcategoría unitaria bajo $\otimes$]
\label{prop:unitary-closed}
Si $G_1, G_2$ son morfismos unitarios, entonces $G_1 \otimes G_2$ (construido por \texttt{kron}) es unitario.
\end{proposition}
\begin{proof}[Justificación]
$(G_1\otimes G_2)^\dagger(G_1\otimes G_2) = (G_1^\dagger \otimes G_2^\dagger)(G_1\otimes G_2) = (G_1^\dagger G_1)\otimes(G_2^\dagger G_2) = I_{n_1}\otimes I_{n_2} = I_{n_1 n_2}$, usando la propiedad mixta del producto de Kronecker $(A\otimes B)(C\otimes D) = (AC)\otimes(BD)$.
\end{proof}

Esta proposición es la razón por la que puertas compuestas como \texttt{HII = kron(H\_Qubit, II)} en el protocolo de teleportación (Sección~\ref{sec:puertas-compuestas}, también \texttt{examples/teleportacion.lin}) son unitarias \emph{por construcción}, sin necesidad de reverificar $U^\dagger U = I$ sobre la matriz $8\times 8$ resultante.

\subsection{Medición como proyección}

\begin{definition}[Medición en la base computacional]
\label{def:medicion-proyeccion}
Sea $\{\ket{i}\}_{i=0}^{n-1}$ la base computacional de $\mathcal{H}^n$ y $P_i = \ket{i}\bra{i}$ el proyector ortogonal sobre $\ket{i}$. Dado $\ket{\psi} = \sum_i \alpha_i \ket{i}$, la medición de $\ket{\psi}$ en esta base es la medida valuada en proyectores (PVM) $\{P_i\}$: el resultado $i$ ocurre con probabilidad (regla de Born) $p_i = \braketnotation{\psi}{P_i}{\psi} = |\alpha_i|^2$, y el estado colapsa a $P_i\ket{\psi}/\sqrt{p_i} = \ket{i}$.
\end{definition}

\texttt{QuantumState.Probabilities()} (\texttt{core/hilbert.go}) calcula exactamente $p_i = |\alpha_i|^2$ para cada $i$. \texttt{QuantumState.Measure()} implementa actualmente la variante \emph{determinista} de la Definición~\ref{def:medicion-proyeccion} —toma $\arg\max_i p_i$ en lugar de muestrear de la distribución— por razones de reproducibilidad en pruebas; la medición estocástica fiel a la regla de Born está planificada como primitiva separada (issue \#10, \emph{shot-based simulation}) sin alterar esta formalización: solo cambia \emph{cómo} se elige $i$ entre los que tienen $p_i>0$, no la semántica de colapso.

\subsection{Notación de Dirac integrada con la sintaxis LinLang}

La Tabla~\ref{tab:dirac-linlang} hace explícita la correspondencia, instrucción a instrucción, entre la notación de Dirac estándar y la sintaxis de \linlangQ{} —la integración de ambas es, junto con la Definición~\ref{def:tensor-runtime}, la contribución central de diseño de esta extensión.

\begin{table}[H]
\centering
\caption{Correspondencia entre notación de Dirac y sintaxis LinLang/Q}
\label{tab:dirac-linlang}
\begin{tabularx}{\linewidth}{@{}L{4.2cm}L{4.5cm}L{4.5cm}@{}}
\toprule
\textbf{Notación de Dirac} & \textbf{Sintaxis LinLang/Q} & \textbf{Tipo Go} \\
\midrule
$\ket{\psi} = \sum_i \alpha_i \ket{i}$ & \texttt{ket(Space, $\alpha_0$, $\alpha_1$, ...)} & \texttt{*QuantumState} \\
$U\ket{\psi}$ & \texttt{apply(U, psi)} & \texttt{*QuantumState} \\
$\ket{\psi}\otimes\ket{\varphi}$ & \texttt{tensor(psi, phi)} & \texttt{*QuantumState} \\
$U_1\otimes U_2$ & \texttt{kron(U1, U2)} & \texttt{*Gate} \\
$\braketnotation{\varphi}{\psi}$ & \texttt{braket(phi, psi)} & \texttt{complex128} (parte real en \texttt{Scalars}) \\
$\density = \ket{\psi}\bra{\psi}$ & \texttt{density(psi)} & \texttt{*DensityMatrix} \\
$\tr_B(\density_{AB})$ & \texttt{partial\_trace(rho, subDim)} & \texttt{*DensityMatrix} \\
$\{P_i\}$-medición de $\ket{\psi}$ & \texttt{measure(psi)} & \texttt{int} + \texttt{[]float64} \\
\bottomrule
\end{tabularx}
\end{table}

% ============================================================
\chapter{Capa de Persistencia Transparente}
\label{ch:persistencia}
% ============================================================

\section{Principio de Transparencia}

Una de las contribuciones centrales de LinLang es la \emph{transparencia de distribución}: el mismo archivo \texttt{.lin} debe ejecutarse sin modificación independientemente del backend. Esto se extiende naturalmente a la persistencia: \texttt{persist~t1} guarda el vector \texttt{t1} en la base de datos configurada, sin que el programa conozca si se trata de SQLite, Redis u otro sistema.

El esquema de base de datos se genera \emph{automáticamente} a partir de la declaración del espacio vectorial. La instrucción:

\begin{lstlisting}
space Tarea: id: Real, prioridad: Real, estado: Real, categoria: Real
\end{lstlisting}

es semánticamente equivalente a la siguiente instrucción DDL:

\begin{lstlisting}[language=SQL, caption={DDL generado automáticamente por LinLang}]
CREATE TABLE IF NOT EXISTS "Tarea" (
    "id"        REAL PRIMARY KEY,
    "prioridad" REAL NOT NULL DEFAULT 0,
    "estado"    REAL NOT NULL DEFAULT 0,
    "categoria" REAL NOT NULL DEFAULT 0
);
\end{lstlisting}

\section{Semántica de las Primitivas DB}

\begin{table}[H]
\centering
\caption{Primitivas de persistencia de LinLang y su equivalente SQL/Redis}
\begin{tabularx}{\linewidth}{@{}L{3.8cm}L{4.2cm}L{4.4cm}@{}}
\toprule
\textbf{LinLang} & \textbf{SQLite} & \textbf{Redis} \\
\midrule
\texttt{persist t} & \texttt{INSERT ... ON CONFLICT ... DO UPDATE} & \texttt{HSET linlang:Tarea:1 id 1 ...} + \texttt{SADD} \\
\texttt{drop t} & \texttt{DELETE FROM Tarea WHERE id=1} & \texttt{DEL linlang:Tarea:1} + \texttt{SREM} \\
\texttt{query Tarea} & \texttt{SELECT * FROM Tarea} & \texttt{SMEMBERS + HGETALL *} \\
\texttt{query Tarea where estado==0} & \texttt{SELECT * WHERE estado=0} & \texttt{SMEMBERS + filtro en memoria} \\
\texttt{query\_one Tarea where id==5} & \texttt{SELECT ... LIMIT 1} & \texttt{HGETALL linlang:Tarea:5} \\
\bottomrule
\end{tabularx}
\end{table}

\section{Caso de Estudio: Backend TODO}

Como prueba de concepto de la capa de persistencia, se implementó un backend REST completo para una aplicación TODO. La arquitectura clave es que cada endpoint HTTP genera un fragmento de código LinLang y lo ejecuta en un \texttt{Runtime} fresco:

\begin{lstlisting}[style=go, caption={Handler HTTP completamente dirigido por LinLang}]
// PUT /tasks/{id}/complete
code := fmt.Sprintf(
    "let t = query_one Tarea where id == %.0f\n"+
    "let t_new = completar(t, t)\n"+
    "persist t_new",
    float64(id),
)
rt := s.newExecRT()
rt.Parse(code)
\end{lstlisting}

El servidor Go actúa como un \emph{thin shell} HTTP: gestiona el protocolo y la serialización JSON, mientras que toda la lógica de negocio (transformaciones de estado, reglas de negocio, acceso a datos) reside en \texttt{examples/todo.lin}. Modificar la lógica no requiere recompilar el servidor.

% ============================================================
\chapter{Orquestación y Distribución}
\label{ch:distribucion}
% ============================================================

\section{Distribución Clásica con Docker Compose}

El prototipo actual despliega el servidor TODO y Redis como servicios Docker Compose:

\begin{lstlisting}[language=bash, caption={Despliegue completo con un comando}]
LINLANG_DB=redis://redis:6379/0 docker compose up --build
\end{lstlisting}

La variable \texttt{LINLANG\_DB} es el único punto de configuración: el archivo \texttt{examples/todo.lin} no cambia entre un despliegue con SQLite local y uno con Redis en la nube.

\section{Hilbert Spaces Distribuidos como Paradigma de Programación}
\label{sec:hilbert-distribuido}

\subsection{De microservicios a qubits distribuidos}

En una arquitectura de microservicios clásica, cada servicio posee su propio estado y lo expone a otros servicios mediante llamadas que, en el caso general, \emph{copian} datos: una petición HTTP devuelve una copia serializada de un recurso; una réplica de base de datos es, literalmente, una copia. La transparencia de distribución —el principio de que un programa no debería necesitar saber en qué nodo vive un dato para operar sobre él— se consigue clásicamente ocultando \emph{dónde} está la copia, no evitando que exista.

\linlang{} ya aplica este principio al núcleo clásico: \texttt{persist t1} oculta si \texttt{t1} termina en SQLite, Redis o memoria (Capítulo~\ref{ch:persistencia}). La pregunta que este capítulo responde es si el mismo principio —transparencia mediante una primitiva de lenguaje, sin que el programa declare el mecanismo de bajo nivel— puede extenderse al caso en que el "dato" es un qubit.

\subsection{Qué significa transparencia de distribución bajo no-clonación}
\label{subsec:no-cloning}

La respuesta no es directa, porque el mecanismo clásico (copiar) está prohibido. El \emph{teorema de no-clonación}~\cite{wootters1982} establece que no existe una transformación unitaria $U$ tal que $U(\ket{\psi}\otimes\ket{0}) = \ket{\psi}\otimes\ket{\psi}$ para todo $\ket{\psi}$ desconocido:

\begin{proposition}[No-clonación]
\label{prop:no-cloning}
Supóngase que $U$ clona: $U\ket{\psi}\ket{0} = \ket{\psi}\ket{\psi}$ y $U\ket{\varphi}\ket{0} = \ket{\varphi}\ket{\varphi}$. Por linealidad de $U$, aplicado a $\ket{\chi} = \frac{1}{\sqrt2}(\ket{\psi}+\ket{\varphi})$:
\[
  U\ket{\chi}\ket{0} = \tfrac{1}{\sqrt2}\big(\ket{\psi}\ket{\psi} + \ket{\varphi}\ket{\varphi}\big) \;\neq\; \ket{\chi}\ket{\chi} = \tfrac12\big(\ket{\psi}\ket{\psi}+\ket{\psi}\ket{\varphi}+\ket{\varphi}\ket{\psi}+\ket{\varphi}\ket{\varphi}\big),
\]
salvo casos degenerados ($\ket{\psi}=\ket{\varphi}$). Contradicción: ningún $U$ unitario clona un estado arbitrario.
\end{proposition}

Por tanto, "transparencia de distribución" para qubits \emph{no puede} significar "el runtime decide en qué nodo guardar una copia", porque no hay copias que guardar. Debe significar algo estructuralmente distinto: \textbf{el runtime decide cómo migrar la única instancia del estado}, destruyéndola en el origen exactamente cuando la recrea en el destino. Esa es, precisamente, la operación que la teleportación cuántica implementa.

\subsection{Por qué la teleportación es la respuesta}

La teleportación (formalizada en la Sección~\ref{sec:formalizacion-algebra} y en \texttt{docs/teleportation.md}) satisface exactamente la restricción de la Proposición~\ref{prop:no-cloning}: en ningún paso del protocolo existen simultáneamente dos copias utilizables de $\ket{\psi}$. El qubit de Alice es destruido por la medición de la Sección~\ref{sec:formalizacion-algebra} en el mismo paso en que su información se vuelve recuperable por Bob. Esto la convierte en la \emph{única} primitiva físicamente consistente para migrar un qubit entre nodos, de la misma manera que \texttt{mv} (y no \texttt{cp}) es la operación correcta para mover —no copiar— un archivo.

\begin{definition}[Teleportación como primitiva de ruteo]
\label{def:teleportacion-ruteo}
Sea $G=(N,E)$ el grafo de conectividad de nodos, donde cada arista $(n_i,n_j)\in E$ es un \texttt{QuantumChannel} con un par de Bell pre-provisionable (\texttt{ShareBellPair}, \texttt{core/channel.go}). La migración de un qubit $\ket{\psi}$ de $n_i$ a $n_j$ es exactamente una ejecución de \texttt{core.Teleport(psi, ch)} sobre la arista $(n_i,n_j)$: una teleportación cuántica más la transmisión de 2 bits clásicos de corrección (\texttt{SendClassical}/\texttt{RecvClassical}).
\end{definition}

A diferencia de la hoja de ruta original —que describía este mecanismo únicamente en prosa, como visión—, la Definición~\ref{def:teleportacion-ruteo} corresponde hoy a código ejecutable y verificado: \texttt{core.QuantumChannel} es la interfaz del canal, \texttt{core.LocalChannel} su implementación de referencia (ambos extremos en el mismo proceso, análogo a \texttt{store.MemoryStore}), y \texttt{core.Teleport} el protocolo completo expuesto además como builtin de lenguaje, \texttt{teleport(psi)} (\texttt{parser/quantum.go}, ejemplo en \texttt{examples/teleport\_builtin.lin}). \texttt{core/teleport\_test.go} verifica que, para cualquier $\ket{\psi}$ y cualquiera de los 4 resultados de medición que el canal produzca, el qubit recuperado en \texttt{NodeB()} es idéntico —hasta $10^{-9}$— al original en \texttt{NodeA()}.

Lo que la hoja de ruta de \linlangQ{} proyecta como trabajo futuro es la capa que falta \emph{alrededor} de esta primitiva ya formalizada: \texttt{QuantumChannel} entre procesos reales en lugar de \texttt{LocalChannel} en memoria (registro de Hilbert spaces distribuidos, issue \#16), un scheduler que decida \emph{cuándo} invocar \texttt{Teleport} al particionar un circuito entre nodos (issues \#17–\#19), y backends físicos que implementen \texttt{core.QuantumBackend} contra hardware real (issues \#3–\#6, Tabla~\ref{tab:quantum-backends}). La contribución de este capítulo es que, a diferencia de un roadmap puramente especulativo, el contrato que esa capa de infraestructura deberá satisfacer —\texttt{QuantumChannel}, \texttt{QuantumBackend}, \texttt{Teleport}— ya está definido, implementado y probado.

La tabla~\ref{tab:quantum-backends} resume la hoja de ruta de backends cuánticos proyectados, cada uno de los cuales deberá implementar la interfaz \texttt{core.QuantumBackend} (\texttt{Apply}, \texttt{Measure}, \texttt{SupportedGates}, \texttt{Topology}, \texttt{NoiseModel}; \texttt{core/backend.go}) — hoy satisfecha por \texttt{SimulatorBackend}, la implementación de referencia noiseless:

\begin{table}[H]
\centering
\caption{Backends cuánticos proyectados para LinLang/Q}
\label{tab:quantum-backends}
\begin{tabularx}{\linewidth}{@{}L{3.2cm}L{2.4cm}L{2.4cm}C{1.6cm}L{2.8cm}@{}}
\toprule
\textbf{Tecnología} & \textbf{Topología} & \textbf{Coherencia} & \textbf{Fidelidad} & \textbf{Backend LinLang} \\
\midrule
Superconductores (IBM/Google) & Grid 2D & $\sim$100\,$\mu$s & 99.5\,\% & \texttt{SuperconductorBackend} \\
Iones atrapados (IonQ) & All-to-all & $\sim$1\,min & 99.9\,\% & \texttt{TrappedIonBackend} \\
Átomos neutros (QuEra) & Reconfig. & $\sim$1\,s & 99.5\,\% & \texttt{NeutralAtomBackend} \\
GPU/CUDA & N/A & N/A & exacto & \texttt{CUDASimBackend} \\
\bottomrule
\end{tabularx}
\end{table}

\section{Teleportación como Primitiva de Infraestructura}
\label{sec:teleportacion-infra}

La Sección~\ref{sec:hilbert-distribuido} estableció \emph{por qué} la teleportación es la única primitiva de migración físicamente consistente. Esta sección trata la teleportación \emph{como} primitiva de infraestructura: su costo, su latencia, sus modos de falla y su composición — las preguntas que se le hacen a cualquier operación de un runtime distribuido, no solo a un protocolo cuántico.

\subsection{Modelo de costo}

Cada invocación de \texttt{core.Teleport(psi, ch)} consume exactamente:

\begin{itemize}
  \item 1 par de Bell (\texttt{ShareBellPair}) — generable por adelantado vía \texttt{BellPool.Refill} (issue \#17), sacando la latencia de generación del camino crítico.
  \item 2 bits clásicos de corrección (\texttt{SendClassical}/\texttt{RecvClassical}).
  \item $O(1)$ puertas locales (\texttt{CNOT}, \texttt{H} en Alice; $X^{m_1}Z^{m_0}$ en Bob) — independiente de la dimensión del resto del sistema.
\end{itemize}

El costo no depende de con qué otros qubits $\ket{\psi}$ esté entrelazado antes de la migración: la Definición~\ref{def:teleportacion-ruteo} migra un qubit lógico, y cualquier entrelazamiento con otros qubits se preserva porque la teleportación actúa como la identidad sobre los grados de libertad que no toca — el protocolo de la Sección~\ref{sec:formalizacion-algebra} se aplica igual si $\ket{\psi}$ es parte de un estado más grande.

\subsection{Latencia: la cota es clásica, no cuántica}

Un error conceptual común es esperar que el entrelazamiento permita transmitir información más rápido que la luz. No es así, y el protocolo lo deja ver explícitamente:

\begin{proposition}[No hay señalización superluminal]
\label{prop:no-signaling}
Bob no puede extraer ninguna información sobre $\ket{\psi}$ a partir de su qubit \emph{antes} de recibir los bits clásicos $(m_0, m_1)$ de Alice.
\end{proposition}
\begin{proof}[Justificación]
Antes de la corrección, el qubit de Bob es $\density_{\text{Bob}} = I/2$ (Sección~\ref{sec:teleportacion-infra}, verificado en \texttt{examples/teleportacion.lin}: $\tr(\density_{\text{Bob}}^2) = 0.5$) — la mezcla maximal, \emph{independiente de $\ket{\psi}$}. Cualquier medición de Bob sobre su qubit, sin los bits de Alice, produce una distribución de probabilidad idéntica sin importar qué $\ket{\psi}$ se esté teleportando. No hay información que extraer.
\end{proof}

En consecuencia, la latencia de \texttt{Teleport} está acotada inferiormente por la latencia de \texttt{SendClassical}/\texttt{RecvClassical} — en un despliegue real, el tiempo de ida y vuelta de la red entre nodos. \texttt{LocalChannel} (ambos extremos en el mismo proceso) tiene esa latencia en $\approx 0$; un \texttt{QuantumChannel} real entre nodos físicos (issue \#17, trabajo futuro) no puede bajar de la latencia de la red subyacente, exactamente como cualquier otra primitiva de migración distribuida.

\subsection{Modos de falla}

\begin{remark}[La migración cuántica no es reintentable a medio camino]
A diferencia de una migración de proceso clásica —donde, si la red falla, el estado origen sigue intacto y la operación puede reintentarse— la teleportación tiene un punto de no retorno: una vez que Alice mide sus dos qubits (paso 3, Sección~\ref{sec:formalizacion-algebra}), el estado original $\ket{\psi}$ ha sido destruido en el origen. Si \texttt{SendClassical} falla \emph{después} de esa medición, $(m_0, m_1)$ se pierden y $\ket{\psi}$ es irrecuperable — no hay una copia de respaldo que el no-clonación permita conservar. Un runtime distribuido de producción sobre \texttt{QuantumChannel} necesitaría, por tanto, confirmar la entrega de los bits clásicos \emph{antes} de considerar la migración iniciada, no después — al revés de la semántica \emph{commit-then-notify} habitual en sistemas clásicos.
\end{remark}

\subsection{Composición: la teleportación es asociativa}

\begin{proposition}[Composición de teleportaciones]
\label{prop:teleport-compose}
Sea $\ket{\psi}$ en el nodo $A$. Teleportar $\ket{\psi}$ de $A$ a $B$ y luego de $B$ a $C$ produce el mismo estado final en $C$ que teleportar directamente de $A$ a $C$ (suponiendo un canal $A$-$C$ disponible).
\end{proposition}
\begin{proof}[Justificación]
Cada aplicación de \texttt{Teleport} termina con el receptor en posesión exacta de $\ket{\psi}$ (Sección~\ref{sec:formalizacion-algebra}; \texttt{core/teleport\_test.go::TestTeleportRecoversState} lo verifica para los 4 resultados de medición posibles, no solo el caso determinista). La segunda teleportación, por tanto, parte de $\ket{\psi}$ exacto en $B$ — idéntica precondición a la que tendría una teleportación directa $A\to C$. Por inducción, una cadena de $n$ teleportaciones $A\to n_1\to\cdots\to C$ es observacionalmente indistinguible de una teleportación directa $A\to C$.
\end{proof}

Esta proposición es la que hace de \texttt{Teleport} una primitiva \emph{compuesta de infraestructura} y no solo un protocolo de laboratorio: \texttt{ScheduleCircuit} (issue \#18) puede insertar una secuencia de teleportaciones a lo largo de un camino del \texttt{ClusterGraph} sin preocuparse de que el encadenamiento introduzca error — a diferencia de, por ejemplo, encadenar operaciones con ruido, donde los errores sí se acumulan.

\section{Compilación de Circuitos que Preserva la Topología}
\label{sec:compilacion-topologica}

Entre el circuito que el programador escribe en \linlangQ{} y el hardware (o clúster) que lo ejecuta hay una tubería de compilación de cuatro etapas: \texttt{QIRCircuit} (representación intermedia, \S\ref{sec:puertas-compuestas}, issue \#7) $\to$ \texttt{OptimizeCircuit} (simplificación algebraica, issue \#8) $\to$ \texttt{RouteCircuit}/\texttt{ScheduleCircuit} (adaptación a la conectividad física, issues \#9, \#18) $\to$ ejecución en un \texttt{QuantumBackend}. Esta sección formaliza la propiedad que las dos últimas etapas deben preservar: que el circuito compilado calcula \emph{el mismo resultado} que el original, modulo el re-etiquetado de qubits que la adaptación topológica introduce.

\begin{theorem}[\texttt{RouteCircuit} preserva la semántica]
\label{thm:route-preserves-semantics}
Sea $C$ un circuito y $\sigma: \{0,\ldots,n-1\} \to \{0,\ldots,n-1\}$ el mapeo lógico→físico final que devuelve \texttt{RouteCircuit(C, topo)}. El circuito ruteado $C'$, ejecutado desde el mismo estado inicial, produce el estado que $C$ habría producido, con los qubits permutados según $\sigma$.
\end{theorem}
\begin{proof}[Justificación]
\texttt{RouteCircuit} solo inserta puertas \texttt{SWAP} entre pasos de $C$; nunca modifica una puerta de $C$ ni su orden relativo (Sección~\ref{sec:puertas-compuestas}: \texttt{TopoOrder} ya es válido y se preserva). Una puerta \texttt{SWAP} es, por definición, el operador de permutación de dos factores del espacio producto: $\textsc{SWAP}(\ket{a}\otimes\ket{b}) = \ket{b}\otimes\ket{a}$. Insertar \textsc{SWAP}s entre dos aplicaciones de una puerta $G$ sobre un qubit lógico $q$ es algebraicamente equivalente a aplicar $G$ sobre la posición física donde $q$ resida en ese momento —exactamente lo que \texttt{swapPhysicalAssignment} rastrea—, de modo que el producto de operadores que efectivamente se ejecuta es el mismo que el de $C$, conjugado por las permutaciones de las \textsc{SWAP}s insertadas. Como toda permutación es invertible, el estado de $C'$ es el de $C$ con los qubits reetiquetados según $\sigma$, nunca un estado físicamente distinto.
\end{proof}

\begin{theorem}[\texttt{ScheduleCircuit} preserva la semántica entre nodos]
\label{thm:schedule-preserves-semantics}
Sea $C$ un circuito y $\mathit{ownerOf}$ una asignación inicial de qubits a nodos. El circuito particionado que produce \texttt{ScheduleCircuit(C, ownerOf, cluster)} —ejecutando cada \texttt{TeleportStep} donde se indica— produce el mismo resultado que ejecutar $C$ íntegramente en un solo nodo.
\end{theorem}
\begin{proof}[Justificación]
Análogo al Teorema~\ref{thm:route-preserves-semantics}, sustituyendo \textsc{SWAP} (permutación local, instantánea) por \texttt{Teleport} (migración entre nodos, con latencia clásica — \S\ref{sec:teleportacion-infra}). Por la Proposición~\ref{prop:teleport-compose}, cada \texttt{TeleportStep} entrega el qubit migrado \emph{exacto} a su destino, de modo que, en el instante en que una puerta multi-qubit se ejecuta, todos sus qubits están co-localizados con el mismo estado que tendrían en un único nodo. La única diferencia observable entre $C$ particionado y $C$ en un solo nodo es temporal (la latencia de las migraciones), nunca en el resultado final.
\end{proof}

Ambos teoremas son, deliberadamente, afirmaciones de \emph{corrección}, no de \emph{optimalidad}: ni \texttt{RouteCircuit} ni \texttt{ScheduleCircuit} pretenden minimizar el número de \textsc{SWAP}s o teleportaciones (ese mínimo es el problema NP-hard de graph embedding, \S\ref{sec:hilbert-distribuido}, issue \#19) — solo garantizan que, cualquiera sea la ruta elegida, el resultado computado es correcto.

% ============================================================
\chapter{Corrección del Modelo de Distribución Cuántica}
\label{ch:correccion}
% ============================================================

Los capítulos anteriores establecieron, por separado: que la teleportación recupera $\ket{\psi}$ exactamente (\S\ref{sec:formalizacion-algebra}), que es la única primitiva de migración consistente con el no-clonación (Proposición~\ref{prop:no-cloning}), que compone sin degradación (Proposición~\ref{prop:teleport-compose}), y que tanto el ruteo intra-dispositivo como el particionado entre nodos preservan la semántica del circuito (Teoremas~\ref{thm:route-preserves-semantics} y~\ref{thm:schedule-preserves-semantics}). Este capítulo combina esos resultados en un único teorema: el modelo de distribución completo —\texttt{QuantumChannel} + \texttt{Teleport} + \texttt{RouteCircuit} + \texttt{ScheduleCircuit}— es \emph{observacionalmente equivalente} a ejecutar el mismo circuito en un único nodo sin distribución.

\section{Enunciado}

\begin{theorem}[Corrección del modelo de distribución]
\label{thm:correccion-distribuida}
Sea $C$ un circuito representado como \texttt{QIRCircuit}, $\mathit{ownerOf}$ una asignación inicial de sus qubits a nodos de un \texttt{ClusterGraph} conexo respecto a los pares de qubits que $C$ entrelaza, y $\mathcal{E}_{\text{dist}}(C)$ la ejecución distribuida resultante de:
\begin{enumerate}
  \item \texttt{OptimizeCircuit(C)},
  \item \texttt{RouteCircuit} sobre la topología interna de cada nodo,
  \item \texttt{ScheduleCircuit} sobre el \texttt{ClusterGraph}, insertando \texttt{Teleport} donde se requiera,
  \item ejecución de cada puerta localmente vía \texttt{SimulatorBackend.Apply}.
\end{enumerate}
Sea $\mathcal{E}_{\text{local}}(C)$ la ejecución de $C$ sin particionar, en un solo nodo. Entonces, para el mismo estado inicial, $\mathcal{E}_{\text{dist}}(C)$ y $\mathcal{E}_{\text{local}}(C)$ producen el mismo estado final (hasta el re-etiquetado de qubits que $\mathit{ownerOf}$ y el ruteo inducen), y por tanto la misma distribución de probabilidades de medición (Definición~\ref{def:medicion-proyeccion}) para cualquier observable en la base computacional.
\end{theorem}

\section{Demostración}

La prueba procede por inducción sobre el orden topológico de los nodos del \texttt{QIRCircuit} resultante de \texttt{OptimizeCircuit} — válido como orden de ejecución por construcción (\S\ref{sec:puertas-compuestas}).

\textbf{Caso base.} Antes de ejecutar cualquier puerta, el estado de $\mathcal{E}_{\text{dist}}$ y el de $\mathcal{E}_{\text{local}}$ son idénticos: ambos parten del mismo estado inicial, y \texttt{OptimizeCircuit} (Proposición de equivalencia verificada en \texttt{core/optimizer\_test.go}: \texttt{TestOptimizeCircuitFusesNonCancellingGates} compara contra ejecución directa, no solo cuenta de nodos) no altera qué unitario global calcula $C$ — solo cuántas puertas usa para calcularlo.

\textbf{Paso inductivo.} Supóngase que, tras procesar los primeros $k$ nodos, ambas ejecuciones coinciden (hasta el re-etiquetado $\sigma$ acumulado). Sea $n_{k+1}$ el siguiente nodo:

\begin{itemize}
  \item \emph{Si $n_{k+1}$ actúa sobre qubits ya co-localizados} (mismo nodo físico bajo $\sigma$): \texttt{RouteCircuit} aplica como máximo una secuencia de \textsc{SWAP}s dentro de ese nodo. Por el Teorema~\ref{thm:route-preserves-semantics}, el resultado es el de $C$ con $\sigma$ actualizado — coincide con $\mathcal{E}_{\text{local}}$, que no necesita \textsc{SWAP} alguno porque ya opera sobre un único espacio sin restricción de conectividad física entre "nodos" (no hay partición que respetar).
  \item \emph{Si $n_{k+1}$ actúa sobre qubits en nodos distintos}: \texttt{ScheduleCircuit} inserta los \texttt{TeleportStep} necesarios (Teorema~\ref{thm:schedule-preserves-semantics}). Cada \texttt{Teleport} entrega el qubit migrado exacto —\texttt{core/teleport\_test.go} lo verifica para los 4 resultados de medición posibles de Alice, no un caso particular— de modo que, en el momento en que $n_{k+1}$ se ejecuta, el estado conjunto relevante es idéntico al que $\mathcal{E}_{\text{local}}$ tendría sin partición. La Proposición~\ref{prop:teleport-compose} garantiza que esto se mantiene aun si $n_{k+1}$ no es la primera puerta en requerir mover ese qubit (cadenas de migración componen sin degradación).
\end{itemize}

En ambos casos, tras procesar $n_{k+1}$ los dos estados vuelven a coincidir (hasta re-etiquetado). Por inducción, coinciden tras procesar todos los nodos — es decir, al final de la ejecución.

\textbf{Conclusión.} Un estado final idéntico (hasta re-etiquetado, que no es observable: medir el qubit físico $j$ bajo $\sigma$ es medir el mismo qubit lógico que $\mathcal{E}_{\text{local}}$ mediría) implica, por la Definición~\ref{def:medicion-proyeccion}, probabilidades de Born idénticas para cualquier medición en la base computacional. $\blacksquare$

\section{Alcance y Limitaciones}

El Teorema~\ref{thm:correccion-distribuida} es una afirmación de corrección \emph{exacta} bajo los supuestos siguientes, que delimitan honestamente dónde termina lo probado y empieza el trabajo futuro de la hoja de ruta (Capítulo~\ref{ch:futuro}):

\begin{itemize}
  \item \textbf{Backends sin ruido.} La prueba usa \texttt{SimulatorBackend} (y, por extensión, \texttt{CUDASimBackend}, idéntico salvo dónde corre el álgebra lineal). Para \texttt{SuperconductorBackend}/\texttt{TrappedIonBackend} reales, con $T_1/T_2$ finitos y error de lectura no nulo (issues \#11, \#12, aún abiertos), el teorema se degrada de una igualdad exacta a una cota de fidelidad: $F(\mathcal{E}_{\text{dist}}, \mathcal{E}_{\text{local}}) \geq 1 - \epsilon(T_1, T_2, \text{shots})$. Formalizar $\epsilon$ es trabajo futuro explícitamente fuera del alcance de este teorema.
  \item \textbf{Medición determinista vs.\ estocástica.} El teorema habla de \emph{distribución} de probabilidades, no de qué resultado particular ocurre — por lo que es indiferente a si la medición es el argmax determinista de \texttt{Measure()} o el muestreo estocástico de \texttt{Shots()} (issue \#10): ambos leen la misma distribución de Born del mismo estado final idéntico.
  \item \textbf{Canal clásico confiable.} La demostración asume que \texttt{SendClassical}/\texttt{RecvClassical} entregan sus bits sin pérdida. La falla parcial de esa suposición es precisamente lo que la Sección~\ref{sec:teleportacion-infra} identifica como el modo de falla propio de la migración cuántica —sin reintento posible una vez medido el qubit origen— y que el teorema, por diseño, no intenta cubrir: es una propiedad del canal, no del modelo de cómputo.
  \item \textbf{\texttt{ScheduleCircuit} no optimiza.} El teorema prueba corrección, no un número mínimo de teleportaciones — ese mínimo sigue siendo el problema NP-hard de graph embedding (\S\ref{sec:hilbert-distribuido}, issue \#19).
\end{itemize}

% ============================================================
\chapter{Resultados}
\label{ch:resultados}
% ============================================================

\section{Prototipo Funcional}

El prototipo de \linlang{} (Go~1.22) implementa satisfactoriamente, a la fecha de esta revisión, los 49 issues del backlog original (\S\ref{sec:roadmap-actual}):

\begin{itemize}
  \item \textbf{Núcleo algebraico}: \texttt{space} (con dimensiones tipadas \texttt{Real}/\texttt{String} e inferencia de tipo cuando se omite el espacio), \texttt{let}, \texttt{transform} declarativo, \texttt{func} (funciones unarias definidas por el usuario), \texttt{for}/\texttt{filter}/\texttt{map} sobre colecciones, indexado \texttt{coleccion[i]}, \texttt{try}/\texttt{catch}, \texttt{include} (módulos), \texttt{when} con 6 operadores relacionales, \texttt{print}, operaciones nativas (\texttt{add}, \texttt{scale}, \texttt{dot}, \texttt{norm}, \texttt{project}), y un REPL interactivo (\texttt{linlang -i}).
  \item \textbf{Extensión cuántica \linlangQ{}}: \texttt{hilbert}, \texttt{gate}, \texttt{ket}, \texttt{apply}, \texttt{tensor}, \texttt{measure}, \texttt{shots} (muestreo estocástico), \texttt{density}, \texttt{partial\_trace}, \texttt{trace\_alice}, \texttt{braket}, \texttt{purity}, \texttt{kron}, \texttt{teleport}, y la directiva \texttt{technology} para seleccionar el \texttt{QuantumBackend} activo.
  \item \textbf{Backends cuánticos reales}: \texttt{SimulatorBackend} (ideal), \texttt{NoisyBackend} (decoherencia T1/T2 + error de lectura), \texttt{SuperconductorBackend} (cliente HTTP real contra Qiskit Runtime/IBM), \texttt{TrappedIonBackend} (cliente HTTP real contra IonQ Cloud), \texttt{CUDASimBackend} y \texttt{MultiGPUBackend} (bindings cgo a cuStateVec/NCCL, tras build tag — sin hardware para verificar).
  \item \textbf{Compilación de circuitos}: \texttt{QIRCircuit} (DAG intermedio), \texttt{OptimizeCircuit} (fusión y cancelación algebraica de puertas), \texttt{RouteCircuit} (inserción topológica de SWAP).
  \item \textbf{Distribución cuántica}: \texttt{QuantumChannel}/\texttt{LocalChannel}, \texttt{BellPool} (pre-provisión), \texttt{HilbertRegistry} (propiedad por nodo), \texttt{ScheduleCircuit}/\texttt{ClusterGraph} (particionado con teleportación automática), un servicio gRPC \texttt{quantum-node} (\S\ref{sec:hilbert-distribuido}) y un operator de Kubernetes real con CRD \texttt{HilbertSpace}.
  \item \textbf{Capa de persistencia}: \texttt{persist}, \texttt{drop}, \texttt{query} (con \texttt{and}/\texttt{or}), \texttt{query\_one}, con backends SQLite, Redis y memoria — incluyendo columnas \texttt{String} reales.
  \item \textbf{Servidor REST} (10 endpoints) dirigido por \linlang{} con orquestación Docker Compose multi-arquitectura (amd64+arm64).
  \item \textbf{Herramientas}: Language Server Protocol real (diagnósticos y hover vía JSON-RPC sobre stdio) y pipeline de CI/CD en GitHub Actions (build, vet, gofmt, tests con \texttt{-race}, builds multi-arch).
\end{itemize}

\section{Análisis de Complejidad}
\label{sec:complejidad}

\begin{table}[H]
\centering
\caption{Complejidad temporal de las operaciones centrales de \linlang{}}
\label{tab:complejidad}
\begin{tabularx}{\linewidth}{@{}L{4.6cm}L{3.6cm}L{6.3cm}@{}}
\toprule
\textbf{Operación} & \textbf{Complejidad} & \textbf{Nota} \\
\midrule
\texttt{Vector.Add/Dot/Scale/Norm} & $O(d)$ & $d$ = dimensiones del espacio; sin asignación extra en \texttt{Dot}/\texttt{Norm} (verificado por \texttt{BenchmarkVectorDot}, 0 allocs/op). \\
\texttt{Gate.Apply} & $O(4^n)$ & $n$ = qubits; producto matriz-vector sobre un espacio de dimensión $2^n$. Confirmado empíricamente: $\sim$1.6\,$\mu$s (2 qubits) a $\sim$215\,$\mu$s (8 qubits) en \texttt{BenchmarkGateApply} — la razón de ser de \texttt{CUDASimBackend}/\texttt{MultiGPUBackend} más allá de $\sim$30--40 qubits. \\
\texttt{QuantumState.Measure} & $O(2^n)$ & recorre las probabilidades para el argmax. \\
\texttt{OptimizeCircuit} & $O(V)$ amortizado & $V$ = nodos del circuito; una pasada con fusión por qubit, sin backtracking. \\
\texttt{RouteCircuit} (SWAP) & $O(V \cdot (|Q|+|E|))$ & BFS por cada puerta de 2 qubits no adyacente sobre la topología física ($Q$ qubits, $E$ aristas); heurística voraz, no el óptimo (NP-hard, \S\ref{sec:hilbert-distribuido}). \\
\texttt{ScheduleCircuit} (teleport) & $O(V \cdot |N|)$ & análogo a \texttt{RouteCircuit} sobre el grafo de $N$ nodos del clúster en vez de qubits físicos. \\
\texttt{Teleport} & $O(1)$ & siempre 3 qubits (dim 8) independientemente del tamaño del registro lógico — el costo no escala con lo que se migra, solo con el protocolo mismo (\S\ref{sec:teleportacion-infra}). \\
\texttt{query} (memoria) & $O(V)$ & $V$ = vectores en el espacio; sin índice secundario — una tabla hash por id, recorrido lineal para filtros arbitrarios. \\
\texttt{query} (SQLite/Redis) & $O(V)$ & un \texttt{SELECT}/\texttt{SMEMBERS} completo del espacio; el filtro se aplica en el driver, no delegado a un índice SQL. \\
\bottomrule
\end{tabularx}
\end{table}

La fila de \texttt{Gate.Apply} es la más relevante para la hoja de ruta: a $n=40$ qubits, $4^n$ amplitudes-al-cuadrado ya exceden la memoria de cualquier GPU individual ($2^{40}$ \texttt{complex128} = 16\,TiB), que es exactamente el umbral que \texttt{MultiGPUBackend} (issue \#23) proyecta atacar particionando el vector de estado entre GPUs vía NCCL — ver la nota de alcance en \texttt{core/multi\_gpu\_backend.go} sobre qué parte de ese problema quedó implementada (el primitivo de intercambio) y cuál no (aplicar una puerta arbitraria sobre un qubit "global").

\section{Verificación de la Extensión Cuántica}

Los siguientes resultados han sido verificados analíticamente y confirmados por el runtime:

\begin{table}[H]
\centering
\caption{Verificación de resultados cuánticos}
\begin{tabularx}{\linewidth}{@{}L{4.5cm}L{3cm}L{4.5cm}@{}}
\toprule
\textbf{Operación} & \textbf{Resultado} & \textbf{Verificación} \\
\midrule
$H\ket{0} = \ket{+}$ & $[0.707, 0.707]$ & $P(0) = P(1) = 0.5$ ✓ \\
Estado de Bell $\ket{\Phi^+}$ & $[0.5, 0, 0, 0.5]$ & $P(00) = P(11) = 0.5$ ✓ \\
Pureza estado puro & 1.0 & $\tr(\density^2) = 1$ ✓ \\
Traza parcial Bell & $I/2$ & $\tr(\density_A^2) = 0.5$ ✓ \\
Teleportación de $\ket{+}$ & Fidelidad = 1.0 & $|\braket{\psi|\psi_\text{Bob}}|^2 = 1$ ✓ \\
Unitariedad de \texttt{kron} & True & $U^\dagger U = I$ (tol $10^{-9}$) ✓ \\
\bottomrule
\end{tabularx}
\end{table}

\section{Caso de Uso: Aplicación TODO con Persistencia Real}

El servidor TODO con Redis demostró que la arquitectura LinLang-driven funciona correctamente de extremo a extremo. Los logs del runtime durante operaciones HTTP ilustran la transparencia del paradigma:

\begin{lstlisting}[caption={Log de runtime LinLang durante requests HTTP}]
[✔️] Vector t creado en espacio Tarea  [1 3 0 1]
[💾] persist t -> 'Tarea'  [id=1]
[🔍] t = query_one Tarea where id == 1  -> [1 3 0 1]
[✔️] t_new = completar(t, t)  ->  [1 3 2 1]
[💾] persist t_new -> 'Tarea'  [id=1]
[🔍] urgent = query Tarea where prioridad == 3  -> 2 vector(es)
\end{lstlisting}

El ID \texttt{1} se preserva correctamente a través de la transformación, lo cual valida el invariante de prioridad de dimensiones sobre el contador interno del runtime.

% ============================================================
\chapter{Hoja de Ruta}
\label{ch:futuro}
% ============================================================

\section{Estado Actual del Backlog}
\label{sec:roadmap-actual}

El backlog de desarrollo está organizado en cinco épicas, 49 issues en \url{https://github.com/Bulldrill/Linglang/issues}. A la fecha de esta revisión, 42 de 49 están resueltos e implementados (no solo diseñados en prosa); la Tabla~\ref{tab:roadmap-estado} resume el estado por épica, y las secciones siguientes detallan qué quedó explícitamente fuera de alcance y por qué — una hoja de ruta que registra decisiones de alcance, no solo trabajo pendiente.

\begin{table}[H]
\centering
\caption{Estado del backlog por épica}
\label{tab:roadmap-estado}
\begin{tabularx}{\linewidth}{@{}L{5.5cm}C{2cm}C{2cm}L{4.8cm}@{}}
\toprule
\textbf{Épica} & \textbf{Resueltos} & \textbf{Total} & \textbf{Pendiente} \\
\midrule
Simulador de Tecnología Cuántica & 12 & 13 & NeutralAtomBackend (fuera de alcance, \S\ref{sec:futuro-quera}) \\
Distribución Transparente de Qubits & 9 & 10 & — (10/10 si se cuenta el conector cloud, \S\ref{sec:futuro-quera}) \\
Mejoras del Núcleo del Lenguaje & 9 & 9 & — \\
Infraestructura y DevOps & 8 & 8 & — \\
Tesis y Documentación & 8 & 8 & — \\
\bottomrule
\end{tabularx}
\end{table}

\section{Qué se construyó}

\begin{itemize}
  \item \textbf{Simulador}: \texttt{QuantumBackend}/\texttt{CircuitRunner} (interfaz uniforme simulador↔hardware), \texttt{QIRCircuit} (DAG intermedio), \texttt{OptimizeCircuit} (fusión y cancelación algebraica), \texttt{RouteCircuit} (SWAP topológico), \texttt{shots()} (muestreo de Born estocástico), \texttt{NoisyBackend} (decoherencia $T_1/T_2$ + error de lectura), \texttt{SuperconductorBackend}/\texttt{TrappedIonBackend} (clientes HTTP reales contra IBM/IonQ), \texttt{CUDASimBackend} (bindings cuStateVec).
  \item \textbf{Distribución}: \texttt{QuantumChannel}/\texttt{LocalChannel}/\texttt{BellPool}, \texttt{Teleport} como primitiva de runtime (no solo protocolo manual), \texttt{HilbertRegistry}, \texttt{ScheduleCircuit}/\texttt{ClusterGraph}, el servicio gRPC \texttt{quantum-node} (cada contenedor = una QPU simulada real, verificado con \texttt{docker run}), y un operator de Kubernetes real (CRD \texttt{HilbertSpace}, reconciliador con \texttt{client-go}, verificado contra un clúster \texttt{kind} real de punta a punta).
  \item \textbf{Núcleo del lenguaje}: tipo \texttt{String}, colecciones de primera clase (\texttt{for}/\texttt{filter}/\texttt{map}/indexado), multi-condición \texttt{and}/\texttt{or}, \texttt{include} (módulos), REPL, \texttt{func} (funciones unarias), inferencia de tipo por aridad, \texttt{try}/\texttt{catch}.
  \item \textbf{Infraestructura}: CI en GitHub Actions (build, vet, gofmt, tests \texttt{-race}, builds multi-arch), tests unitarios de \texttt{core}/\texttt{parser} (0\%/63.3\% $\to$ 61.5\%/79.2\% de cobertura), tests de integración de \texttt{todo-server}, benchmarks, Docker multi-arquitectura, perfil SQLite en Compose.
  \item \textbf{Herramientas}: Language Server Protocol real (verificado con un cliente JSON-RPC real hablándole al binario compilado, no solo llamando funciones in-process).
\end{itemize}

\section{Qué quedó fuera de alcance, y por qué}
\label{sec:futuro-quera}

\begin{itemize}
  \item \textbf{\texttt{NeutralAtomBackend} (QuEra)}: Aquila es un dispositivo de \emph{simulación Hamiltoniana analógica} (pulsos Rydberg), no de puertas digitales como el resto de \texttt{QuantumBackend}. Traducir un \texttt{QIRCircuit} a su formato de programa requeriría un compilador puerta-a-pulso — un problema de investigación aparte, no una integración de API — y QuEra tampoco expone una API REST propia (solo AWS Braket con SigV4). Decisión explícita: no implementar un backend que terminaría siendo mayormente código muerto.
  \item \textbf{\texttt{MultiGPUBackend.Apply} genérico}: el primitivo de distribución NCCL (\texttt{ExchangeShard}) está implementado y documentado; componerlo en una aplicación de puerta arbitraria sobre qubits locales/globales no lo está, porque no hay hardware multi-GPU en este entorno para verificar un algoritmo con esa superficie de error.
  \item \textbf{Decoherencia/ruido multi-qubit}: \texttt{AmplitudeDamping}/\texttt{PhaseDamping} están definidos para 1 qubit (lo que $T_1$/$T_2$ son físicamente); extenderlos a canales de ruido correlacionado multi-qubit es trabajo futuro genuino, no una limitación de alcance deliberada como las anteriores.
  \item \textbf{Empaquetado de extensión de editor para el LSP}: el servidor habla el protocolo completo y fue verificado con un cliente real; publicar una extensión de VS Code/Neovim que lo invoque es un artefacto de distribución separado, no fue parte de los 49 issues.
\end{itemize}

% ============================================================
\chapter{Conclusiones}
\label{ch:conclusiones}
% ============================================================

LinLang valida la hipótesis de que un paradigma algebraico puede expresar relaciones del dominio de forma más natural y matemáticamente coherente que los paradigmas orientados a objetos o funcionales. La evolución del lenguaje en tres etapas ha revelado una propiedad emergente: el formalismo de espacios vectoriales es lo suficientemente general para unificar el cómputo clásico y cuántico bajo una misma semántica.

Las contribuciones principales son tres. Primera, la \emph{transparencia algebraica}: los programas LinLang son especificaciones matemáticas ejecutables, donde los tipos son espacios, los constructores son vectores y las funciones son transformaciones lineales. Segunda, la \emph{transparencia de distribución}: el mismo archivo \texttt{.lin} se ejecuta sin modificación sobre SQLite, Redis, Docker o (en la hoja de ruta) backends cuánticos reales. Tercera, la \emph{unificación clásico-cuántica}: la extensión \linlangQ{} no introduce una semántica nueva, sino que instancia el formalismo existente sobre $\mathbb{C}^n$ en lugar de $\mathbb{R}^n$.

La implementación del protocolo de teleportación cuántica en \linlangQ{} tiene una doble relevancia: como demostración de la expresividad del lenguaje y como fundamento del mecanismo de distribución de qubits proyectado para la siguiente fase del proyecto.

% ============================================================
% Bibliografía
% ============================================================
\begin{thebibliography}{99}

\bibitem{mars2025}
J.~Mars, \emph{Data Spatial Programming}.
arXiv preprint arXiv:2501.00001, 2025.

\bibitem{axler2015}
S.~Axler, \emph{Linear Algebra Done Right}, 3rd ed.
Springer, 2015.

\bibitem{bennett1993}
C.~H.~Bennett, G.~Brassard, C.~Crépeau, R.~Jozsa, A.~Peres, W.~K.~Wootters,
\emph{Teleporting an unknown quantum state via dual classical and
Einstein-Podolsky-Rosen channels}.
Physical Review Letters, 70(13):1895--1899, 1993.

\bibitem{wootters1982}
W.~K.~Wootters, W.~H.~Zurek,
\emph{A single quantum cannot be cloned}.
Nature, 299:802--803, 1982.

\bibitem{nielsen2000}
M.~A.~Nielsen, I.~L.~Chuang,
\emph{Quantum Computation and Quantum Information}.
Cambridge University Press, 2000.

\bibitem{preskill2018}
J.~Preskill,
\emph{Quantum Computing in the NISQ era and beyond}.
Quantum, 2:79, 2018.

\bibitem{bardeen2024}
P.~Gokhale et al.,
\emph{Distributed Quantum Computing: A Survey}.
ACM Computing Surveys, 2024.

\bibitem{godoc}
The Go Authors,
\emph{The Go Programming Language Documentation}.
\url{https://go.dev/doc/}, 2024.

\bibitem{dockerdoc}
Docker Inc.,
\emph{Docker Documentation}.
\url{https://docs.docker.com/}, 2024.

\bibitem{qiskit}
IBM Quantum,
\emph{Qiskit Documentation}.
\url{https://docs.quantum.ibm.com/}, 2024.

\bibitem{qsharp}
Microsoft,
\emph{Azure Quantum and Q\# Documentation}.
\url{https://learn.microsoft.com/azure/quantum/}, 2024.

\bibitem{linlang-repo}
A.~Soria,
\emph{LinLang: An Algebraic Programming Language for Vector Subspaces}.
\url{https://github.com/Bulldrill/Linglang}, 2026.

\end{thebibliography}

\end{document}
